Claims that a model \emph{cannot} do something (impossibility,
information limits, irreducible error) are treated as durable
conclusions, yet they are as prone to artifact as claims that a model
\emph{can}. We document a worked example end to end. Investigating the
canonical ``should I eat this mushroom?'' decision, we built a careful,
reproducible, adversarially-verified study and concluded that the task
is \textbf{information-limited}: on real image embeddings a strong
aggregate separability (AUROC 0.94) coexisted with an apparently
irreducible \textasciitilde5\% confidently-wrong ``false-eat'' residual
that \emph{no method in a panel of \textasciitilde18} could remove, with
formal divergence-floor and data-processing arguments to explain why,
and decision-curve analysis showing abstention dominates. The conclusion
survived an adversarial-verification process that had already caught two
of our own \emph{positive} results as rigged. It was wrong anyway.
Within the same pipeline and the same pairs, a stronger frozen encoder
(DINOv2) cut the false-eat to near zero: \textbf{≤0.7\%} per cell
un-cropped (≤1.0\% under a substrate crop; pooled 95\% upper bound
\textbf{0.4\%}), where the weak encoder left \textbf{0.1--21.3\%}, lower
within every pair. This included the freshly-scraped textbook
``you-need-a-microscope'' confusions (honey-mushroom vs deadly
\emph{Galerina}, where it cleared 0 of 150 held-out poisons). What
matters is the encoder-dependence, not the absolute rate. (The
strong-encoder cells are near-zero rather than exact zeros: at the
original ≈60-specimen sampling they read a bare ``0.0\%'' whose ≈4.9\%
upper bound overlapped the original residual. At full sampling they read
0.1--0.7\%, with a single-cell bound ≈2\% and a pooled bound ≈0.4\%.) An
information limit is a property of the \emph{photograph} and would bind
every encoder equally, so a stronger backbone that beats a weaker one
within every pair is incompatible with one: the ``information limit''
was a \textbf{weak-baseline artifact}. These random-sample rates,
however low, do not bound the worst-case individual specimen, which we
could not test and which remains the thesis's one honest residual (§6).
We show this across \textbf{six lethal look-alike decisions spanning
five biological groups} (fungi, plants, reptiles, arachnids,
amphibians), \textbf{three anchored to externally-published claims} (two
of visual indistinguishability, one of systematic misidentification:
coral-snake mimicry, \citep{pfennig2001}; brown-recluse
misidentification, \citep{vetter2000}; poison-frog mimicry,
\citep{symula2001}). In every group the weak encoder leaves a large
false-safe rate (often near-chance, clearing up to half of the dangerous
specimens) that a stronger frozen encoder beats within every pair: to
near zero in five of six groups, and always far below the weak baseline.
One small-\emph{n} poison-frog pair is the honest exception, at
\textasciitilde10\% against the weak encoder's 33\%. We offer the
representation-invariance test as a portable protocol for
frozen-representation limit claims, not a general theorem about all
safety-negatives, and we distill the post-mortem (a weak baseline, a
curated benchmark mistaken for the field, and a
true-but-empirically-vacuous theory whose antecedent we never checked)
into a checklist. The lesson generalizes beyond mushrooms:
\emph{benchmark separability is not field safety, and a limitation claim
is only as strong as the baseline that fails to beat it.}

\begin{center}\rule{0.5\linewidth}{0.5pt}\end{center}
